% Insert after the all-exponential-orders theorem and its proof.
% Uses the notation and coefficient formula of angular_transseries.tex.
\subsection{The complete exponential series diverges at every fixed index}
\label{angular:sec:divergence}

The finite-cutoff theorem cannot be strengthened to ordinary convergence
of the complete coefficient series. This obstruction persists after
all equal rational exponential bases have been combined.
For a base $q$, put
\[
 \widehat C(q;b)=\sum_{\mathbf m:\,q(\mathbf m)=q}C_{\mathbf m}(b).
\]
Each such sum is finite by local finiteness of the scale.

\begin{theorem}[Divergence after grouping equal bases]
\label{angular:thm:divergence}
For every fixed integer $a\ge1$ and real $b>0$, the ordinary series
\[
 \sum_q\widehat C(q;b)q^a
\]
does not converge. In fact its individual terms are unbounded along
an infinite sequence of distinct rational bases tending to zero.
The ungrouped multiset series also diverges. Thus the complete
expansion in Theorem~\ref{angular:thm:allorders} is asymptotic and
necessarily divergent at every fixed $a,b$.
\end{theorem}
\begin{proof}
Choose an integer $s\ge1$ with $2s>a$, and put $K=2s+1$.
For an odd prime $P\ne3$, consider the multiset with $2s$ copies
of $3$ and one copy of $P$. Its base is
\begin{equation}\label{angular:eq:prime-base}
 q_P=\frac{2^K}{3^{2s}P}.
\end{equation}
We first prove that this base has no other multiset representation.
If $k$ integers $n_1,\ldots,n_k\ge3$ give the same base, then
\[
 \prod_{j=1}^k n_j=2^{k-K}3^{2s}P.
\]
Since the left side is an integer and $3^{2s}P$ is odd, $k\ge K$.
Let $\Omega(n)$ denote the number of prime factors of $n$, with
multiplicity. Additivity gives
\[
 \sum_{j=1}^k\Omega(n_j)=(k-K)+2s+1=k.
\]
Each summand is at least one, so every $n_j$ is prime. The lower
bound $n_j\ge3$ excludes the prime two; hence $k=K$, and unique
factorization leaves precisely $2s$ copies of $3$ and one copy
of $P$. Therefore $\widehat C(q_P;b)$ is exactly the coefficient
of this one multiset, with no possible cancellation from other
representations.

Write $h_j=H_j^{(b)}$, let $\sigma_P=\sin(P\pi/2)\in\{-1,1\}$,
and introduce the even analytic germ
\[
 B(u)=\frac{u}{\sin(2u)},\qquad B(0)=\frac12.
\]
For odd $P$, the coefficient germs in
\eqref{angular:eq:phi} satisfy
\[
 \phi_3(u)=B(u)\cos(3u),\qquad
 \phi_P(u)=-\sigma_P B(u)\cos(Pu).
\]
The multinomial prefactor in \eqref{angular:eq:coefficient} is
$(K-1)!/((2s)!1!)=1$. Consequently
\begin{align}
 \widehat C(q_P;b)
 &=-\sigma_Ph_2^{2s}h_{P-1}\,\mathcal P_s(P),
 \label{angular:eq:prime-coefficient}\\
 \mathcal P_s(P)
 &:=[u^{2s}]B(u)^{2s+1}\cos^{2s}(3u)\cos(Pu)
 \notag\\
 &=\frac{(-1)^s}{2^{2s+1}(2s)!}P^{2s}
     +O_s(P^{2s-2}).
 \label{angular:eq:prime-polynomial}
\end{align}
The last equality is a polynomial identity in $P$ followed by its
leading-term estimate. Its highest power comes only from the
$u^{2s}$ coefficient of $\cos(Pu)$ and the constant coefficient
$2^{-2s-1}$ of the remaining factor. All other powers are even
and at most $2s-2$.

Combining \eqref{angular:eq:prime-base} and
\eqref{angular:eq:prime-polynomial} gives, along odd primes,
\[
 \bigl|\widehat C(q_P;b)q_P^a\bigr|
 =\frac{2^{K(a-1)}h_2^{2s}}{3^{2sa}(2s)!}
   h_{P-1}P^{2s-a}\bigl(1+O_s(P^{-2})\bigr).
\]
Here $h_{P-1}\ge1$. Since $2s-a>0$, these terms tend to infinity
as $P\to\infty$ through odd primes different from three. Such
primes are infinite; no estimate for their distribution is needed.
The bases $q_P$ are distinct and tend to zero. The necessary
condition that individual summands tend to zero therefore fails,
both for the grouped series and for the ungrouped series, and
indeed for every ordering of these ordinary series.
\end{proof}

For example, the first two polynomials in
\eqref{angular:eq:prime-polynomial} are
\[
 \mathcal P_1(P)=-\frac{P^2+14}{16},\qquad
 \mathcal P_2(P)=\frac{P^4}{768}+\frac{11P^2}{48}
                   +\frac{81}{32}.
\]
The supplied script \texttt{check\_angular\_divergence.py} computes
these coefficient polynomials exactly through $s=6$, checks their
leading coefficients, and compares six small cases with direct
symbolic trigonometric expansion. The general proof above, not
these finite checks, establishes divergence.

There is no conflict with the complete asymptotic theorem. Its
cutoff is fixed first and then $a$ tends to infinity. The divergent
subsequence above fixes $a$ first and lets the retained frequency
grow without bound. Summation procedures beyond ordinary series
convergence are a separate question.
